\documentclass[a4paper,11pt]{amsart}
\usepackage[margin=1in]{geometry}
\usepackage{amsmath,amssymb,amsthm,mathtools}
\usepackage[T1]{fontenc}
\usepackage{lmodern}
\usepackage{microtype}
\usepackage{enumitem}
\usepackage[hidelinks]{hyperref}

\newtheorem{theorem}{Theorem}[section]
\newtheorem{lemma}[theorem]{Lemma}
\newtheorem{proposition}[theorem]{Proposition}
\newtheorem{corollary}[theorem]{Corollary}
\theoremstyle{remark}
\newtheorem{remark}[theorem]{Remark}

\DeclareMathOperator{\Norm}{N}
\newcommand{\Z}{\mathbb Z}
\newcommand{\Zi}{\mathbb Z[i]}

\title[A Ljunggren--Nagell Reduction]{A Ljunggren--Nagell Reduction}
\author{}
\date{}

\begin{document}

\begin{abstract}
We prove that the Diophantine equation
\[
        y(x^3-z^2)=x
\]
has no solution in positive integers.  More precisely, we prove that the only
solution with $x,y\geq 1$ and $z\geq 0$ is $(x,y,z)=(1,1,0)$.  The elementary
part of the argument reduces the problem to the classical Ljunggren--Nagell
input that the equation $U^2+1=A^4B^3$ has no positive solution.  All
reductions after that input are given in full detail.
\end{abstract}

\maketitle

\section{The external input}

We use the following standard Ljunggren--Nagell theorem.  It is the only
external ingredient in the proof.

\begin{theorem}[Ljunggren--Nagell input]\label{thm:LN}
The equation
\[
        U^2+1=A^4B^3
\]
has no solution in integers $U,A,B$ with $U\geq 1$ and $A,B\geq 1$.  Equivalently,
the only solution with $U\geq0$ and $A,B\geq1$ is
\[
        (U,A,B)=(0,1,1).
\]
\end{theorem}

\begin{remark}
The rest of the manuscript is self-contained once Theorem~\ref{thm:LN} is
accepted.  Appendix~\ref{app:gaussian} records the Gaussian-integer form of
Theorem~\ref{thm:LN}, because it explains exactly why this is the natural
classical input for the present equation.
\end{remark}

\section{Elementary reductions}

We shall repeatedly use the following elementary fact.

\begin{lemma}[Coprime factors of a square]\label{lem:coprime-square}
Let $r,s\in\Z_{>0}$ satisfy $\gcd(r,s)=1$.  If $rs$ is a square, then both
$r$ and $s$ are squares.
\end{lemma}

\begin{proof}
Write prime factorizations
\[
        r=\prod_p p^{\alpha_p},\qquad s=\prod_p p^{\beta_p},
\]
where all but finitely many exponents are zero.  The condition
$\gcd(r,s)=1$ says that for each prime $p$, at most one of
$\alpha_p,\beta_p$ is nonzero.  Since $rs$ is a square, each
$\alpha_p+\beta_p$ is even.  Therefore the nonzero exponent among
$\alpha_p,\beta_p$, if there is one, is even.  Thus every prime exponent in
$r$ and in $s$ is even, so both $r$ and $s$ are squares.
\end{proof}

\begin{lemma}[The quotient forced by the equation]\label{lem:quotient}
Let $x,y,z\in\Z$ satisfy
\[
        x\geq1,
        \qquad y\geq1,
        \qquad z\geq0,
        \qquad y(x^3-z^2)=x.
\]
Then $y\mid x$.  If $d=x/y$, then
\[
        z^2=d(d^2y^3-1).                     \tag{2.1}\label{eq:z2-d}
\]
Moreover, if $z>0$, then $d^2y^3-1>0$.
\end{lemma}

\begin{proof}
The defining equation gives
\[
        x^3-z^2=\frac{x}{y}.
\]
The left-hand side is an integer.  Hence $x/y$ is an integer, so $y\mid x$.
Put
\[
        d=\frac{x}{y}\in\Z_{>0}.
\]
Then $x=dy$, and substitution gives
\[
        y\bigl((dy)^3-z^2\bigr)=dy.
\]
Dividing by $y>0$ yields
\[
        d^3y^3-z^2=d,
\]
so
\[
        z^2=d^3y^3-d=d(d^2y^3-1),
\]
which is \eqref{eq:z2-d}.  If $z>0$, then $z^2>0$, and since $d>0$, the
factor $d^2y^3-1$ must also be positive.
\end{proof}

\begin{lemma}[The square quotient reduction]\label{lem:square-quotient}
Let $x,y,z\in\Z$ satisfy
\[
        x\geq1,
        \qquad y\geq1,
        \qquad z\geq0,
        \qquad y(x^3-z^2)=x.
\]
Then either $(x,y,z)=(1,1,0)$, or there are integers $a,u\geq1$ such that
\[
        x=a^2y,
        \qquad z=au,
        \qquad u^2+1=a^4y^3.                 \tag{2.2}\label{eq:LN-form}
\]
\end{lemma}

\begin{proof}
By Lemma~\ref{lem:quotient}, $d=x/y$ is a positive integer and
\[
        z^2=d(d^2y^3-1).
\]
First suppose
\[
        d^2y^3-1=0.
\]
Then $d^2y^3=1$.  Since $d,y\in\Z_{>0}$, this forces $d=1$ and $y=1$.
Thus $x=dy=1$, and the displayed formula for $z^2$ gives $z=0$.  Hence
$(x,y,z)=(1,1,0)$.

It remains to consider the case
\[
        d^2y^3-1>0.
\]
The two positive factors $d$ and $d^2y^3-1$ are coprime, because any common
divisor divides both $d$ and $d^2y^3-1$, hence divides
\[
        d^2y^3-(d^2y^3-1)=1.
\]
Their product is the square $z^2$.  Lemma~\ref{lem:coprime-square} therefore
implies that both factors are squares.  Hence there are $a,u\in\Z_{>0}$ such
that
\[
        d=a^2,
        \qquad d^2y^3-1=u^2.
\]
The first equality gives
\[
        x=dy=a^2y.
\]
The second gives
\[
        u^2+1=d^2y^3=a^4y^3,
\]
which is \eqref{eq:LN-form}.  Finally,
\[
        z^2=d(d^2y^3-1)=a^2u^2,
\]
and since $z,a,u\geq0$, with $a,u>0$, we get $z=au$.
\end{proof}

\section{The solution of the original equation}

\begin{theorem}\label{thm:main-nonnegative}
The only integer solution of
\[
        y(x^3-z^2)=x
\]
with $x,y\geq1$ and $z\geq0$ is
\[
        (x,y,z)=(1,1,0).
\]
\end{theorem}

\begin{proof}
Let $x,y,z$ be a solution with $x,y\geq1$ and $z\geq0$.  By
Lemma~\ref{lem:square-quotient}, either $(x,y,z)=(1,1,0)$, or there are
integers $a,u\geq1$ such that
\[
        u^2+1=a^4y^3.
\]
The latter is impossible by Theorem~\ref{thm:LN}, applied with
\[
        U=u,
        \qquad A=a,
        \qquad B=y.
\]
Therefore only the first alternative can occur, and the only solution is
$(1,1,0)$.
\end{proof}

\begin{corollary}\label{cor:positive}
The equation
\[
        y(x^3-z^2)=x
\]
has no solution in positive integers $x,y,z$.
\end{corollary}

\begin{proof}
A positive-integer solution would in particular be a solution with
$x,y\geq1$ and $z\geq0$.  Theorem~\ref{thm:main-nonnegative} says that the
only such solution is $(1,1,0)$, but this has $z=0$, not $z>0$.  Hence no
positive-integer solution exists.
\end{proof}

\begin{remark}
The theorem above proves nonexistence in the positive range and identifies the
single boundary solution with $z=0$.  It does not claim a classification when
$x,y,z$ are allowed to be arbitrary signed integers.
\end{remark}

\appendix

\section{Gaussian-integer form of the quoted input}\label{app:gaussian}

This appendix is included to make the role of the classical input transparent.
It is not a replacement for Theorem~\ref{thm:LN}; rather, it shows why the
input is exactly the Ljunggren--Nagell obstruction naturally produced by the
present Diophantine equation.

Recall that $\Zi$ is a unique factorization domain with units
$\{\pm1,\pm i\}$ and norm
\[
        \Norm(r+si)=r^2+s^2.
\]

\begin{proposition}\label{prop:gaussian-reformulation}
Let $U,A,B\in\Z$ satisfy $U\geq1$ and $A,B\geq1$.  If
\[
        U^2+1=A^4B^3,                         \tag{A.1}\label{eq:LN-app}
\]
then $U$ is even, $A$ and $B$ are odd, and there exist Gaussian integers
$\alpha,\beta\in\Zi$ and a unit $\varepsilon\in\{\pm1,\pm i\}$ such that
\[
        U+i=\varepsilon\alpha^4\beta^3,
        \qquad \Norm(\alpha)=A,
        \qquad \Norm(\beta)=B.                \tag{A.2}\label{eq:gaussian-form}
\]
Consequently Theorem~\ref{thm:LN} is equivalent to the assertion that no
nonunit Gaussian data in \eqref{eq:gaussian-form} can have imaginary part
$1$.
\end{proposition}

\begin{proof}
First consider parity.  If $U$ is odd, then
\[
        U^2+1\equiv 2\pmod 4.
\]
The right-hand side $A^4B^3$ is either odd, if $A$ and $B$ are both odd, or
is divisible by $8$ if $B$ is even, or by $16$ if $A$ is even.  It is never
congruent to $2\pmod4$.  Thus $U$ is even.  Then $U^2+1$ is odd, so
$A^4B^3$ is odd, and therefore $A$ and $B$ are odd.

In $\Zi$ we have
\[
        (U+i)(U-i)=U^2+1=A^4B^3.              \tag{A.3}\label{eq:gauss-factor}
\]
We claim that $U+i$ and $U-i$ are coprime in $\Zi$.  Indeed, any common
divisor divides their difference $2i$.  Hence the norm of any common divisor
is a divisor of $\Norm(2i)=4$.  On the other hand, the norm of any divisor of
$U+i$ divides $\Norm(U+i)=U^2+1$, which is odd.  Therefore the norm of a
common divisor is $1$, so the common divisor is a unit.

Since $\Zi$ is a unique factorization domain and the two factors in
\eqref{eq:gauss-factor} are coprime, every Gaussian prime occurring in
$U+i$ occurs there with exactly the exponent with which its rational norm
contributes to $A^4B^3$.  Thus the prime factorization of $U+i$ can be grouped
into a fourth-power part, a third-power part, and a unit.  Concretely, there
are $\alpha,\beta\in\Zi$ and a unit $\varepsilon$ such that
\[
        U+i=\varepsilon\alpha^4\beta^3.
\]
Taking norms gives
\[
        U^2+1=\Norm(\alpha)^4\Norm(\beta)^3.
\]
The grouping may be chosen so that
\[
        \Norm(\alpha)=A,
        \qquad
        \Norm(\beta)=B,
\]
because the rational prime exponents belonging to $A$ are placed in
$\alpha$, while those belonging to $B$ are placed in $\beta$.

Finally, $U+i$ has imaginary part $1$, so \eqref{eq:gaussian-form} says that
a Gaussian product of a fourth power and a third power has imaginary part
$1$.  Conversely, any identity of the form \eqref{eq:gaussian-form} with
imaginary part $1$ gives, after taking norms, an equation of the form
\eqref{eq:LN-app}.  This proves the claimed equivalence.
\end{proof}

\end{document}
